% Generated draft: requires mathematical and source review.
\documentclass[11pt]{article}
\usepackage[margin=25mm]{geometry}
\usepackage{fontspec}
\setmainfont{Latin Modern Roman}
\usepackage{amsmath,amssymb,mathtools,unicode-math}
\setmathfont{Latin Modern Math}
\usepackage{enumitem,needspace,xurl,hyperref}
\setlength{\parindent}{0pt}
\setlength{\parskip}{5pt}
\setlength{\emergencystretch}{3em}
\newcommand{\N}{\mathbb{N}}
\newcommand{\Z}{\mathbb{Z}}
\newcommand{\Q}{\mathbb{Q}}
\newcommand{\R}{\mathbb{R}}
\newcommand{\C}{\mathbb{C}}
\newcommand{\F}{\mathbb{F}}
\newcommand{\A}{\mathbb{A}}
\newcommand{\PP}{\mathbb{P}}
\newcommand{\E}{\mathbb{E}}
\newcommand{\eps}{\varepsilon}
\newcommand{\dd}{\,\mathrm{d}}
\begin{document}
\section*{UnsolvedMath ID: 6800024}
\textbf{Source problem code:} catalogue entry
\subsection{Definitions and mathematical statement}
The requested lower bound is
\[
d\leq\operatorname{Index}(\Phi_{\mathcal A})+\operatorname{Null}(\Phi_{\mathcal A}),
\]
where the source identifies $\operatorname{Null}(\Phi_{\mathcal A})$ as the dimension of the space of Jacobi fields. The supplied summary interprets $\Phi_{\mathcal A}$ as a minimax minimal surface produced by a $d$-dimensional admissible family $\mathcal A$.

For a smooth closed immersed minimal surface in a fixed Riemannian ambient manifold, a usual geometric convention takes the area second-variation form $Q$ on normal variations. Its Morse index is the maximal dimension of a subspace on which $Q$ is negative definite, and nullity is the dimension of the kernel of the associated normal Jacobi operator. With these conventions the bound says that negative and zero second-variation directions together account for at least $d$ parameters. It is not an upper bound on index and is not replaced by one.

The record does not define $d$, the admissible minmax class, its topology or nontriviality, the construction of $\Phi_{\mathcal A}$, or the ambient and domain categories. Parameter dimension alone is insufficient without an essentiality condition: redundant parameters cannot automatically force additional Jacobi or negative directions. It is also unclear whether the conclusion concerns every minimax surface or existence of a suitable choice.

The Jacobi-field convention must specify normal variations versus reparametrization directions, boundary conditions if any, and treatment of branched, multiple-covered or multiplicity-weighted limits. The standard smooth closed normal-variation description above is explicit context, not a claim that these hypotheses were supplied. The source fragment therefore fixes the exact intended inequality but not all data needed for a self-contained minimax theorem.
\subsection{Short English statement}
Show that the essential parameter dimension of the source's minimax family is no greater than the resulting minimal surface's Morse index plus nullity. The minmax class, output-selection and Jacobi-field conventions are missing.
\subsection{Sources}
\begin{itemize}
\item Local UnsolvedMath registry, permanent ID 6800024, exact source code catalogue entry, ``Minimax minimal surfaces''; proposed by Tristan Rivi\`ere, ETH Z\"urich, 2018. Registry status open, classification OPEN-TRIAGE.
\item Morgan and Pansu, \emph{A List of Open Problems in Differential Geometry} (2018), question environment 24, \url{https://www.imo.universite-paris-saclay.fr/~pansu/problems_MTDG.tex}. Supplied source-TeX extraction accessed July 29, 2026.
\item The supplied summary interprets $d$ as the family dimension in Rivi\`ere's minimax programme and reports no confirmed proof of this precise inequality. It gives no exact later reference or missing definition of the admissible class.
\item Provenance flags current status and redistribution terms NEEDS\_REVIEW; source TeX is publicly available on a coauthor's publication page.
\end{itemize}
\end{document}
